\section{Preliminaries}\label{sec:prelim}
\subsection{Structural Causal Models and the Standard Fairness Model}
\label{sec:sfm}

A structural causal model (SCM) $\mathcal{M} = (V, U, \mathcal{F}, P_U)$
consists of endogenous variables $V$, exogenous background variables $U$,
structural equations~$\mathcal{F}$, and a distribution $P_U$ over exogenous
noise \citep{pearl2000causality}.
The associated causal directed acyclic graph (DAG) $\mathcal{G}$ encodes the
functional dependencies among variables in~$V$.

\noindent
Counterfactual fairness frameworks built on a fully specified SCM can provide
strong, individual-level guarantees, but the assumptions behind them are
restrictive and often untestable. Rather than positing a full SCM, we use the
\textbf{Standard Fairness Model} (SFM) of \citet{plecko2022causalfairnessanalysis},
which projects the observed features onto a compact causal structure that
captures only the roles relevant to fairness analysis.

\begin{definition}[Standard Fairness Model Projection]
\label{def:sfm}
We partition the observed feature space as follows:
\begin{itemize}[nosep]
  \item $X\in\{x^{(0)},x^{(1)}\}$: a binary sensitive attribute (e.g., sex), with $x^{(0)}$ the disadvantaged group;
  \item $W$: mediators, i.e., attributes causally downstream of $X$ that precede the outcome (e.g., education, occupation, hours worked);
  \item $Z$: pre-treatment covariates that are not causally downstream of $X$ (e.g., birthplace, age);
  \item $Y$: the binary outcome, used to train a classifier $\hat{f}$.
\end{itemize}
\end{definition}

A fixed classifier $\hat f:\mathcal X\times\mathcal W\times\mathcal Z\to[0,1]$
produces the score $\hat Y=\hat f(X,W,Z)$ and a positive decision when
$\hat Y\ge\tau$. \emph{Recourse candidates} are disadvantaged-group true
negatives: individuals $(x^{(0)},w_0,z)$ with $\hat f(x^{(0)},w_0,z)<\tau$ and
$y=0$. The induced graph $\mathcal G_{\mathrm{SFM}}$ is complete across
clusters. It contains the direct path $X\to\hat Y$, the indirect paths
$X\to W\to\hat{Y}$ (including pathways within $W$), and the spurious paths
$X\leftrightarrow Z\to\hat{Y}$ and $X\leftrightarrow Z\to W\to \hat{Y}$.
Throughout, the object of study is the prediction $\hat Y$, not $Y$: recourse
acts on the decision rule, so every effect below is an effect \emph{on the
classifier's score}.

\paragraph{Identification of classifier effects.}\label{sec:identification-assumptions}
For $s\in\{x^{(0)},x^{(1)}\}$, let $W_s$ denote the potential mediator, the
value $W$ would take had $X$ been set to $s$. We assume:
\begin{itemize}[nosep]
    \item[(i)] $Z$ contains no descendant of $X$;
    \item[(ii)] \emph{consistency}: $W = W_s$ whenever $X=s$;
    \item[(iii)] \emph{positivity}: $\Pr[X=s \mid z] > 0$ for all $z$ and $s$;
    \item[(iv)] \emph{conditional ignorability}: $W_s \perp X \mid Z$.
\end{itemize}
Assumption~(i) makes it safe to condition on $Z$: holding $z$ fixed blocks the
spurious $X$--$Z$ association, so it is not counted as mediated disparity.
Consistency~(ii) links observed and potential mediators. Positivity~(iii)
guarantees that $P(W\mid X=s,z)$ is defined for both groups at every
recipient's covariates. Ignorability~(iv) rules out unmeasured $X$--$W$
confounding. Together they give
\begin{equation}
  P(W_s\mid z)=P(W\mid X=s,z).
  \label{eq:mediator_id}
\end{equation}
No assumption about confounding of $Y$ is needed. Because $\hat f$ is a known
function, the step $W\to\hat Y$ cannot be confounded. Reading the quantities
below as effects on the true outcome $Y$ would additionally require the
standard sequential-ignorability conditions, which we do not claim.
%We further make the assumption that attributes in $W$ are \emph{actionable}, meaning they can be feasibly modified by the individual (e.g., credit utilization), and hence is a valid mediating pathway.


\subsection{Mediation Analysis on the Prediction Scale}
\label{sec:mediation_background}

We adapt the natural direct and indirect effects
\citep{robins1992identifiability,pearl2001direct} to the classifier's score.
For $a,b\in\{x^{(0)},x^{(1)}\}$, abbreviated $0,1$ in subscripts, define
\begin{equation}
  \mu_{ab}(z)\;:=\;\mathbb E\bigl[\hat f(a,W_b,z)\mid Z=z\bigr],
  \label{eq:mu_ab}
\end{equation}
where $a$ is the value of $X$ that enters $\hat f$ directly and $b$ is the value
that generates the mediators.

\begin{definition}[Natural Direct and Indirect Effects on the prediction]
\label{def:nde_nie}
The \textbf{natural direct effect} changes $X$ in $\hat f$ while holding the
mediators at the values they would naturally take under $x^{(0)}$; the
\textbf{natural indirect effect} holds $X=x^{(0)}$ in $\hat f$ and changes only
the law generating the mediators:
\begin{equation}
  \NDE(z)=\mu_{10}(z)-\mu_{00}(z),\qquad
  \NIE(z)=\mu_{01}(z)-\mu_{00}(z).
  \label{eq:effects}
\end{equation}
\end{definition}

\noindent Both are \emph{pure} effects, measured from the same baseline
$\mu_{00}(z)$, the disadvantaged group's own prediction. Each isolates one
pathway: the $\NDE$ moves only the direct input $X\to\hat Y$, and the $\NIE$
moves only the mediator law. For a nonlinear classifier they need not sum to
the total effect, and we do not use the total effect.

\paragraph{Identification.}
By~\eqref{eq:mediator_id}, each $\mu_{ab}(z)$ is an expectation of a known
function under an observational mediator law,
\begin{equation}
  \mu_{ab}(z)=\mathbb E_{W\sim P(W\mid X=b,\,z)}\bigl[\hat f(a,W,z)\bigr],
  \label{eq:mu_id}
\end{equation}
which is Pearl's mediation formula with $\mathbb E[Y\mid x,w,z]$ replaced by
$\hat f(x,w,z)$. Two features of the prediction scale simplify the classical
setting. First, since $\hat f(a,\cdot,z)$ is deterministic, $\mu_{ab}(z)$
depends on $W_b$ only through its conditional law, so the natural effects
coincide with their interventional (randomized-mediator)
analogues~\citep{didelez2006direct,vanderweele2014effect} and no cross-world
independence assumption is needed. Second, every effect is conditional on $z$,
so it can be evaluated at each recourse candidate's own covariates. The
conditional mediator laws $P(W\mid X,Z)$ are estimated with a block-ordered
normalizing flow (Section~\ref{sec:flows}).

\noindent Because recourse cannot change $X$, it can at most close the mediated
part of the disparity, the $\NIE$. The $\NDE$, carried along $X\to\hat Y$, is
out of reach of any action on $W$. Comparing the two describes how much of the
disparity recourse could address in principle. It does not say how much is
actionable, since not every mediator can be changed.


\subsection{From the Indirect Effect to a Post-Recourse Disparity}
\label{sec:mediated_contrast}

The NIE compares two \emph{natural} mediator laws. Recourse replaces one of
them with an \emph{intervened} law, so our evaluation needs a slightly more
general quantity. For fixed $z$ and any distribution $Q$ over $\mathcal W$, let
\begin{equation}
  \psi_z(Q)\;:=\;\mathbb E_{W\sim Q}\bigl[\hat f(x^{(0)},W,z)\bigr]
  \qquad\text{and}\qquad
  \Delta_z(Q)\;:=\;\psi_z\bigl(P_z^{(1)}\bigr)-\psi_z(Q),
  \label{eq:delta}
\end{equation}
with $P_z^{(s)}:=P(W\mid X=s,z)$. The reference term
$\psi_z(P_z^{(1)})=\mu_{01}(z)$ is fixed: it is the expected score a
disadvantaged individual with covariates $z$ would receive if their mediators
were distributed as in the advantaged group. Only the comparison law $Q$
varies. The NIE is the case in which $Q$ is the disadvantaged group's own
natural law, $\NIE(z)=\Delta_z(P_z^{(0)})$.

For a recourse candidate $i$ with factual mediators $w_{0,i}$, we use three
comparison laws:
\begin{align}
  D_{\mathrm{pre},i}&=\Delta_{z_i}\bigl(P_{z_i}^{(0)}\bigr)=\NIE(z_i),
  &&\text{natural disadvantaged law,}\label{eq:dpre}\\
  D^{\mathrm{fact}}_{\mathrm{pre},i}&=\Delta_{z_i}\bigl(\delta_{w_{0,i}}\bigr),
  &&\text{candidate's factual mediators,}\label{eq:dpre_fact}\\
  D_{\mathrm{post},i}&=\Delta_{z_i}\bigl(Q_{a_i}\bigr),
  &&\text{mediators after action plan } a_i,\label{eq:dpost}
\end{align}
where $Q_{a_i}$ is the interventional law of $W$ induced by the candidate's
action (Section~\ref{sec:intervention}). Only~\eqref{eq:dpre} is a natural
indirect effect. $Q_{a_i}$ is not the law of any potential mediator $W_s$: it
clamps the acted mediators, keeps the candidate's upstream values, and
regenerates the downstream ones. $D_{\mathrm{post},i}$ therefore contrasts a
natural mediator law with an intervened one, in the spirit of
stochastic-intervention effects~\citep{diaz2012population}, and we call it the
\emph{post-recourse mediated disparity} rather than an NIE. Since all three
share the same reference term,
\begin{equation}
  D_{\mathrm{post},i}
  \;=\;D^{\mathrm{fact}}_{\mathrm{pre},i}
  \;-\;\underbrace{\bigl[\psi_{z_i}(Q_{a_i})-\hat f(x^{(0)},w_{0,i},z_i)\bigr]}_{
     \text{effect of the action on the candidate's score}}.
  \label{eq:dpost_decomp}
\end{equation}
Recourse reduces the mediated disparity by exactly the expected score gain it
produces, and $\mu_{01}(z_i)$ is the target that gain should reach. A negative
$D_{\mathrm{post},i}$ is overshoot: the action raised the expected score above
what the advantaged mediator law would give.

Candidates are selected on their factual features, through
$\hat f(x^{(0)},w_{0,i},z_i)<\tau$, so $w_{0,i}$ is not a representative draw
from $P_{z_i}^{(0)}$. We therefore use the NIE decomposition~\eqref{eq:effects}
to describe the disparity at the population level, and measure closure
against $D^{\mathrm{fact}}_{\mathrm{pre},i}$, the baseline that
\eqref{eq:dpost_decomp} compares recourse with.
